\documentclass[10pt,a4paper]{amsart}
\usepackage[margin=19mm]{geometry}
\usepackage{amsmath,amssymb,mathtools}
\usepackage[scr=rsfs,cal=euler]{mathalfa}
\usepackage{xfrac}
\usepackage[unicode,hidelinks,bookmarks=false]{hyperref}
\newcommand{\ie}{i.e.\ }
\newcommand{\cf}{cf.\ }
\newcommand{\resp}{respectively\ }
\newcommand{\Subsection}{\subsection}
\providecommand{\nobreakdash}{\nobreak\hbox{-}}
\setlength{\emergencystretch}{2em}
\multlinegap0pt
\usepackage{amssymb}
\usepackage{enumerate}
\usepackage{bm}
\newtheorem{lemma}{Lemma}
\DeclareMathOperator{\Div}{div}
\DeclareMathOperator*{\argmin}{arg\,min}
\newcommand{\dx}{\,dx}
\title{Global convergence of $W^{1,\infty}$-steepest descent for PDE constrained shape optimisation with semilinear elliptic equations in function space}
\author{Klaus Deckelnick, Philip J. Herbert, Michael Hinze}
\date{Source: arXiv:2603.02812v1 (CC BY 4.0)}
\begin{document}
\maketitle

\noindent\emph{Source and licence.} Global convergence of $W^{1,\infty}$-steepest descent for PDE constrained shape optimisation with semilinear elliptic equations in function space. Version arXiv:2603.02812v1 (CC BY 4.0), distributed by arXiv under the Creative Commons Attribution 4.0 International licence, http://creativecommons.org/licenses/by/4.0/, as recorded in the arXiv licence metadata for this exact version and shown on the abstract page https://arxiv.org/abs/2603.02812v1. Full licence notice: \texttt{licenses/arxiv-2603.02812-CC-BY-4.0.txt}.\par\smallskip

\noindent\textit{Editorial scope.} Counted unit: the article's lemma decrease (source0.tex lines 438--448). The article's complete original proof of it is delivered with it (source0.tex lines 449--629, one \texttt{proof} environment of 181 lines). No mathematics is written, repaired or re-derived: the statement and the proof are the article's own lines, in the article's own notation, and no retained line is substituted or re-flowed.  Retained uncounted as the source-local context the proof needs: lines 231--354; lines 272--274; lines 280--282.  Nothing was omitted from any retained span, so the package carries no omission record.  The retained lines cite 3 works, whose entries are transcribed verbatim from the article's own bibliography source in the e-print and delivered with short editorial labels recorded in the item's reference mapping.  No retained line contains a figure environment, an image-inclusion command or drawing code, so no image is delivered and the package declares no asset dependency.  Editorial formatting only, in addition: an AMS \texttt{amsart} preamble that restates the article's own declarations and macros used by the retained lines, namely the environments \texttt{lemma} on separate counters, together with the standard packages those lines need.  The retained environments print in this document as Lemma 1 in order of appearance, whereas the article prints the same items with the numbering of its own full document.  The complete native source of the article as distributed in the arXiv e-print for this exact version is bundled unchanged as \texttt{sources/195640/source0.tex}.
\noindent
Let $D \subset \mathbb R^d$ be an open, convex, and bounded hold-all domain with a Lipschitz boundary and 
\begin{displaymath}
\mathcal S:= \lbrace \Omega \subset D \, | \, \Omega \mbox{ is open } \rbrace.
\end{displaymath}
\noindent
We consider the shape optimisation problem
\begin{displaymath}
\min_{\Omega \in \mathcal S} \mathcal J(\Omega) = \int_{\Omega} j(\cdot,u,\nabla u)  \dx,
\end{displaymath}
where $u \in H^1_0(\Omega)\cap L^\infty(\Omega)$ is the unique solution of 
\begin{equation} \label{state}
\int_{\Omega} \nabla u \cdot \nabla \eta  \dx + \int_{\Omega} g(u) \eta \dx  = \int_{\Omega }f \eta \dx  \qquad \mbox{ for all } \eta \in H^1_0(\Omega).
\end{equation}
Here  '$\cdot$' denotes the Euclidean inner product of two vectors. 
In what follows we assume that $f \in H^1(D)$, $g \in C^2(\mathbb{R})$ with $g'(t) \ge 0$ for all $t \in \mathbb R$ and that $j \in C^2(D \times \mathbb R \times \mathbb R^d)$ satisfies 
\begin{eqnarray}
| j(x,u,z) | + | j_x(x,u,z)| + | j_{xx}(x,u,z)|  & \leq  & \varphi_1(x)+ c_1 \bigl( | u|^q + |z|^2 \bigr); \label{jest} \\
| j_u(x,u,z) | + | j_{xu}(x,u,z) | & \leq  &  \varphi_2(x) + c_2 \bigl(  |u |^{q-1} + | z|^{2-\frac{2}{q}} \bigr);   \label{j1est} \\
| j_z(x,u,z) | + | j_{xz}(x,u,z) |  & \leq   & \varphi_3(x) + c_3 \bigl( |u|^{\frac{q}{2}} + |z| \bigr);  \label{j2est}   \\
| j_{uu}(x,u,z) | & \leq  &\varphi_4(x) + c_4 \bigl( |u|^{q-2} + |z|^{2- \frac{4}{q}} \bigr); \label{j3est} \\
  | j_{zz}(x,u,z) | & \leq & \varphi_5(x) \label{j4est}
\end{eqnarray}
for all $(x,u,z) \in  D \times  \mathbb R \times \mathbb R^d$, {where a subscript $x$ denotes the derivative with respect to the first variable, $u$, the second, and $z$, the third}.
Here, $2 \leq q < \infty$ if $d=2$, and $q= 6$ if $d =3$. Also, $\varphi_1,\ldots,\varphi_5$ are non-negative with 
$\varphi_1 \in L^1(D), \varphi_2 \in L^{\frac{q}{q-1}}(D), \varphi_3 \in L^2(D), \varphi_4 \in L^{\frac{q}{q-2}}(D)$ and 
$\varphi_5 \in L^\infty(D)$.
{Note that the choice of $q$ implies the continuous embedding $H^1_0(D) \hookrightarrow L^q(D)$, so that there exists $c>0$ with
\begin{equation} \label{Dembed}
\Vert v \Vert_{L^q} \leq c \Vert v \Vert_{H^1} \qquad \mbox{ for all } v \in H^1_0(D).
\end{equation} } 
\noindent
From here onwards we omit the explicit dependencies on the spatial variable $x$ wherever appropriate. 
Using the techniques described in Sections 4.4 and 4.5 of \cite{ADJ21}, one calculates for the shape  derivative of $\mathcal J$
\begin{eqnarray}
 \mathcal J'(\Omega)[V] 
 & = &  \int_{\Omega} \Bigl( j(\cdot,u,\nabla u) \Div V + j_x(\cdot,u,\nabla u)\cdot  V  - j_z(\cdot,u,\nabla u) \cdot DV^{\mathsf{T}} \nabla u \Bigr)  \dx    \label{firstvar} \\
 & & + \int_{\Omega} \Bigl(  \bigl( DV + D V^{\mathsf{T}} - \Div V I \bigr)  \nabla u \cdot \nabla p - g(u)p \mbox{ div }V + {\Div}( f V) p   \Bigr)  \dx \nonumber
\end{eqnarray}
for all $V \in W^{1,\infty}_0(D, \mathbb R^d)$. Here,
$p \in H^1_0(\Omega)$ is the solution of the adjoint problem
\begin{equation} \label{adj}
\int_{\Omega} \nabla p \cdot \nabla \eta  \dx + \int_{\Omega} g'(u) p \eta \dx = \int_{\Omega} \bigl( j_u(\cdot,u,\nabla u) \eta + j_z(\cdot,u,\nabla u) \cdot \nabla \eta \bigr)  \dx \quad \mbox{ for all } \eta \in H^1_0(\Omega).
\end{equation}
Note that the assumptions \eqref{jest}--\eqref{j2est} ensure that
all integrals that appear on the right hand side of \eqref{firstvar} and \eqref{adj} exist. At the same time these growth conditions
also guarantee uniform control on the solutions of the state
and adjoint equations as shown in 
\begin{lemma} \label{lem:apriori} Let $\Omega \in \mathcal S$. Then the problems \eqref{state} and \eqref{adj} admit unique solutions $u\in H^1_0(\Omega)\cap L^\infty(\Omega)$, and $p \in H^1_0(\Omega)$, respectively. Moreover, there exists a constant $c^\star$, such that
\begin{equation} \label{discapriori}
\Vert u \Vert_{H^1} + \Vert u\Vert_{L^\infty} + \Vert p \Vert_{H^1} \leq c^\star.
\end{equation}
The constant $c^*$ only depends on $d,D, f,j$ and $g$, but
is independent of $\Omega$. Furthermore, we think of $u$ and $p$ as being extended by zero to $D$.
\end{lemma}
\begin{proof} In order to prove the existence of a solution $u \in H^1_0(\Omega)$ of
\eqref{state} we first consider for $M>0$ the modified nonlinearity
\begin{displaymath}
g_M(t):= \left\{
\begin{array}{cl}
g(-M), & t<-M, \\
g(t), & -M \leq t \leq M, \\
g(M), & t>M.
\end{array}
\right.
\end{displaymath}
Note that $g_M$ is increasing since $g' \geq 0$. 
Using the theory of monotone operators it can be shown that the problem
\begin{displaymath}
\int_\Omega \nabla u_M \cdot \nabla \eta \dx + \int_\Omega g_M(u_M) \eta \dx =
\int_\Omega f \eta \dx \qquad \forall \eta \in H^1_0(\Omega)
\end{displaymath}
has a unique solution $u_M \in H^1_0(\Omega)$, see \cite[Section 2.6]{VM25} for details.  Let us next denote by
$v \in H^2(D) \cap H^1_0(D)$ the solution to the Poisson problem
\begin{displaymath}
- \Delta v = f - g(0) \; \mbox{ a.e. in }  D, \quad v=0 \mbox{ on }  \partial D.  
\end{displaymath}
Since $H^2(D) \hookrightarrow C^0(\bar D)$ for $d=2,3$ we infer that there exists $K \geq 0$, which only depends on $d,D,f$ and $g$,
such that $| v(x) | \leq K$ for all $x \in \bar D$. The function $w:=v+K$ then satisfies
$w \geq 0$ in $\bar D$ as well as
\begin{displaymath}
- \Delta w + g_M(w) = -\Delta v + g_M(w) = f - g_M(0) + g_M(w) \geq f \quad \mbox{ a.e. in } D
\end{displaymath}
since $g_M$ is increasing and $g_M(0)=g(0)$. If we multiply the above relation by
$\eta \in C^1_0(\Omega), \eta \geq 0$ in  $\Omega$ and use an approximation argument we obtain
\begin{displaymath}
    \int_\Omega \nabla w \cdot \nabla \eta \dx + \int_\Omega g_M(w) \eta \dx \geq
    \int_\Omega f \eta \dx \qquad \forall \eta \in H^1_0(\Omega), \, \eta \geq 0 \mbox{ a.e. in } \Omega.
\end{displaymath}
Hence,
\begin{displaymath}
    \int_\Omega \nabla( u_M-w) \cdot \nabla \eta \dx + \int_\Omega (g_M(u_M)-g_M(w)) \eta \dx \leq
0 \qquad \forall \eta \in H^1_0(\Omega), \eta \geq 0 \mbox{ a.e. in } \Omega.
\end{displaymath}
Since $w \geq 0$ in $D$ we have that  $\eta =(u_M-w)^+ \in H^1_0(\Omega), \eta \geq 0$ and
therefore by the monotonicity of $g_M$
\begin{displaymath}
    \int_\Omega | \nabla(u_M-w)^+ |^2 \dx \leq 0.
\end{displaymath}
Thus, $(u_M - w)^+=0$ so that $u_M \leq w =v+K \leq 2K$ a.e. in $\Omega$. A similar argument
shows that $u_M \geq -2K$ a.e. in $\Omega$, so that 
\begin{equation} \label{eq:umbound}
\Vert  u_M \Vert_{L^\infty}  \leq M
\end{equation}
if we choose $M=2K$. In particular, $g_M(u_M)=g(u_M)$ and $u=u_M$ is a solution of
\eqref{state}, while uniqueness is a consequence of the monotonicity of $g$. While the $L^\infty$--bound for $u$ in \eqref{discapriori} then follows from \eqref{eq:umbound}, the $H^1$--seminorm is controlled
by testing \eqref{state} with $\eta=u$ and using that $g(u)u
\geq g(0)u$ a.e. in $\Omega$. \\
In order to prove the bound on $p$ we test \eqref{adj} with $\eta=p \in H^1_0(\Omega)$ and use 
the fact that $g' \geq 0$, \eqref{j1est}, \eqref{j2est},
H\"older's inequality and  {\eqref{Dembed}} to obtain
\begin{eqnarray*}
\lefteqn{  \int_D | \nabla p |^2  \dx = 
\int_{\Omega} | \nabla p |^2  \dx    } \\
 & \leq & \int_{\Omega} \left[ \bigl( \varphi_2 + c_2( | u |^{q-1} + | \nabla u |^{\frac{2(q-1)}{q}} ) \bigr) | p | + \bigl( \varphi_3 + c_3( | u|^{\frac{q}{2}} + | \nabla u | ) \bigr)  | \nabla p | 
 \right] \dx \\
& \leq & \bigl( \Vert \varphi_2 \Vert_{L^{\frac{q}{q-1}}} + c_2(\Vert u \Vert_{L^q}^{q-1} + \Vert \nabla u \Vert_{L^2}^{\frac{2(q-1)}{q}} ) \bigr)  \Vert p \Vert_{L^q}   \\
& & + \bigl( \Vert \varphi_3 \Vert_{L^2} + c_3( \Vert u \Vert_{L^q}^{\frac{q}{2}} + \Vert \nabla u \Vert_{L^2})  \bigr) \Vert  \nabla p \Vert_{L^2} \\[2mm]
& \leq & c \bigl( 1+ \Vert u \Vert_{H^1}^{q-1}  \bigr)   \Vert p \Vert_{H^1} \leq c \bigl( 1 + (c^*)^{q-1} \bigr) \Vert p \Vert_{H^1} \leq c \bigl( 1 + (c^*)^{q-1} \bigr) \Vert \nabla p \Vert_{L^2},
\end{eqnarray*}
{where we also made use of Poincar\'e's inequality for $D$ and the bound on $u$.} The estimate for 
$\Vert p \Vert_{L^2}$ now follows from another application  of Poincar\'e's inequality.
\end{proof}
\noindent

\begin{equation} 
\int_{\Omega} \nabla p \cdot \nabla \eta  \dx + \int_{\Omega} g'(u) p \eta \dx = \int_{\Omega} \bigl( j_u(\cdot,u,\nabla u) \eta + j_z(\cdot,u,\nabla u) \cdot \nabla \eta \bigr)  \dx \quad \mbox{ for all } \eta \in H^1_0(\Omega).
\end{equation}

\begin{equation} 
\Vert u \Vert_{H^1} + \Vert u\Vert_{L^\infty} + \Vert p \Vert_{H^1} \leq c^\star.
\end{equation}

\begin{lemma} \label{decrease}
 Let {$\gamma \in (0,1)$ be a fixed constant,} $\Omega \in \mathcal S$, and
\begin{displaymath}
V \in \argmin \lbrace J'(\Omega)[W] \, | \, W \in W^{1,\infty}_0(D, \mathbb R^d), \, | DW | \leq 1 \mbox{ a.e. in } D \rbrace.
\end{displaymath}
Suppose that $\mathcal J'(\Omega)[V] \leq - \epsilon$ for some $\epsilon>0$. Then there exists 
$0< \delta <1$ which only depends on $j, f, g, D, d, \gamma$ and $\epsilon$ such that
\begin{displaymath}
\mathcal J(\Omega_{t}) - \mathcal J(\Omega) \leq \gamma t \mathcal J'(\Omega)[V] \qquad \mbox{ for all } 0 \leq t \leq \delta, \; \mbox{ where } \Omega_{t} = T_t(\Omega) \mbox{ and } T_t=\mbox{id}+ t V.
\end{displaymath}
\end{lemma}

\begin{proof}
The proof follows the lines of the proof of \cite[Lemma 3.2]{DHH25-convergence}, in which
a corresponding result was obtained in a finite--dimensional setting, where
state, adjoint state and deformation field are elements of finite element spaces. 
Nevertheless many arguments immediately carry over to the infinite--dimensional case and
we will direct the reader to the corresponding argument in \cite{DHH25-convergence}
when appropriate. \\
Let $\Omega_t=T_t(\Omega)$ with $T_t(x)=x+t V(x)$. In view of the definition of $\mathcal J$
we have
\begin{displaymath}
\mathcal J(\Omega_t) = \int_{\Omega_t} j(\cdot,u_t,\nabla u_t)  \dx,
\end{displaymath}
where $u_t \in H^1_0(\Omega_t)$ solves
\begin{displaymath}  
\int_{\Omega_t} \nabla u_t \cdot \nabla \eta_t  \dx + \int_{\Omega_t} g(u_t) \eta_t \dx = \int_{\Omega_t} f \eta_t  \dx \qquad \forall \eta_t \in H^1_0(\Omega_t).
\end{displaymath}
Choosing $\eta_t:= \eta \circ T_t^{-1} \in H^1_0(\Omega_t)$ for a fixed $\eta \in H^1_0(\Omega)$ we thus infer that 
\begin{equation}  \label{hatuht}
\int_{\Omega_t} \nabla u_t \cdot \nabla (\eta \circ T_t^{-1})  \dx + \int_{\Omega_t} g(u_t) \eta \circ T_t^{-1} \dx = \int_{\Omega_t} f \eta \circ T_t^{-1}  \dx \qquad
\forall  \eta \in H^1_0(\Omega).
\end{equation}
Transforming to $\Omega$ we obtain
\begin{equation} \label{uhtrel}
\int_{\Omega} \nabla u_t \circ T_t \cdot \nabla (\eta \circ T_t^{-1}) \circ T_t \, | \mbox{det} D T_t |  \dx + \int_{\Omega} g(u_t\circ T_t) \eta \, | \mbox{det} D T_t |dx =
\int_{\Omega} f \circ T_t \, \eta \, | \mbox{det} D T_t |  \dx 
\end{equation}
for all $\eta \in H^1_0(\Omega)$. Since $\mbox{det} DT_t = \mbox{det}(I+ t DV)$ and $| D V | \leq 1$ in $\bar D$ one easily
verifies that 
\begin{equation} \label{dif1}
\mbox{det} DT_t -1 = t {\Div} V + r_1, \quad \mbox{ with } | r_1 | \leq c t^2,
\end{equation}
where the constant $c$ only depends on $d$. This also implies that  there exists $\delta_1>0$ such
that $\mbox{det} DT_t >0, 0 \leq t \leq \delta_1$. Defining $A_t:= (DT_t)^{-1} (D T_t)^{- \mathsf{T}} \mbox{det} DT_t$ we hence see that the function
$\hat u_t:= u_t \circ T_t \in H^1_0(\Omega)$ satisfies
\begin{equation} \label{stateuht}
\int_{\Omega} A_t \nabla \hat u_t \cdot \nabla \eta  \dx + \int_{\Omega} g(\hat u_t) \eta \, \mbox{det} D T_t \dx = \int_{\Omega} f \circ T_t \, \eta \, \mbox{det} D T_t  \dx \qquad
\mbox{ for all } \eta \in H^1_0(\Omega).
\end{equation}
As a result we deduce that
\begin{eqnarray}
\mathcal J(\Omega_t) - \mathcal J(\Omega) 
& = & \int_{\Omega} \bigl( j(T_t,\hat u_t, D T_t^{- \mathsf{T}} \nabla \hat u_t) \, \mbox{det} DT_t - j(\cdot,u,\nabla u) \bigr)  \dx   \nonumber \\
& = & \int_{\Omega} j(\cdot,u,\nabla u) (\mbox{det} DT_t -1 )  \dx + \int_{\Omega} \bigl( j(T_t,\hat u_t,DT_t^{- \mathsf{T}} \nabla \hat u_t) - j(\cdot,u,\nabla u) \bigr)  \dx    \nonumber  \\
& & +  \int_{\Omega} \bigl( j(T_t,\hat u_t, DT_t^{- \mathsf{T}} \nabla \hat u_t) -j(\cdot,u,\nabla u) \bigr) ( \mbox{det} DT_t - 1)  \dx  \nonumber \\
& = &  \sum_{j=1}^3 T_j.  \label{difj}
\end{eqnarray}
Arguing as in \cite[(3.7)]{DHH25-convergence} and using Lemma \ref{lem:apriori} we obtain
\begin{equation} \label{eq:t1}
T_1   \leq t \int_{\Omega} j(\cdot,u,\nabla u) \Div V  \dx + c t^2.
\end{equation} 
The term $T_2$ is written in the same way as in \cite[(3.8)]{DHH25-convergence} so that
\begin{eqnarray}
T_2 & = & t  \int_{\Omega} j_x(\cdot,u,\nabla u) \cdot V  \dx - t \int_{\Omega} j_z(\cdot,u,\nabla u) \cdot DV^{\mathsf{T}} \nabla u \nonumber \\
& & + \int_{\Omega} j_z(\cdot,u,\nabla u) \cdot \bigl( (DT_t^{-\mathsf{T}} -I + t DV^{\mathsf{T}} ) \nabla \hat u_t + t DV^{\mathsf{T}} \nabla (u - \hat u_t) \bigr)  \dx \nonumber \\
& & +  \int_{\Omega} \bigl( j_u(\cdot,u,\nabla u)(\hat u_t -u)  + j_z(\cdot,u,\nabla u) \cdot \nabla (\hat u_t - u) \bigr)  \dx \nonumber \\
& & + \int_{\Omega}  \int_0^1 (1-s) \frac{d^2}{ds^2} \left[ j(\cdot +stV,s \hat u_t +(1-s) u,
s  DT_t^{-\mathsf{T}} \nabla \hat u_t +(1-s) \nabla u) \right] ds  \dx \nonumber \\
& = & \sum_{j=1}^5 T_{2,j}. \label{t2}
\end{eqnarray}
The integrals $T_{2,3}$ and $T_{2,5}$ are handled in the same way as in \cite[(3.9), (3.17)]{DHH25-convergence} and therefore
\begin{equation} \label{eq:t23}
T_{2,3} +T_{2,5}  \leq c \bigl(  t^2 + c \Vert \hat u_t - u \Vert_{H^1}^2 \bigr). 
\end{equation} 
In order to treat $T_{2,4}$ we use  \eqref{adj}, \eqref{state} and \eqref{stateuht} and write
\begin{eqnarray*}
T_{2,4} &=& \int_{\Omega} \nabla p \cdot \nabla ( \hat u_t - u )  \dx + \int_{\Omega} g'(u) p (\hat u_t-u) \dx = \int_{\Omega} \nabla p \cdot \nabla \hat u_t  \dx +  \int_{\Omega}g(\hat u_t)p \dx \\
& & - \int_{\Omega} \nabla p \cdot \nabla u  \dx  - \int_{\Omega}g(u)p \dx 
 + \int_{\Omega} \bigl( g'(u)( \hat u_t -u) - g(\hat u_t) + g(u)) p \dx  \\
 & = & \int_{\Omega} A_t \nabla \hat u_t \cdot \nabla p \dx + \int_{\Omega} g(\hat u_t) p \, 
 \mbox{det} D T_t \dx - \int_{\Omega} \nabla p \cdot \nabla u  \dx  - \int_{\Omega}g(u)p \dx \\
 & & + \int_{\Omega} (I -A_t) \nabla \hat u_t \cdot \nabla p \dx + \int_{\Omega} g(\hat u_t) 
 (1 - \mbox{det} D T_t) p \dx \\
 & & + \int_{\Omega} \bigl( g'(u)( \hat u_t -u) - g(\hat u_t) + g(u)) p \dx  \\
 & = & \int_{\Omega}  \bigl( f \circ T_t \, \mbox{det} D T_t -f \bigr) p \dx
+ \int_{\Omega} (I -  A_t) \nabla \hat u_t \cdot \nabla p \dx + \int_{\Omega} g(\hat u_t) 
 (1 - \mbox{det} D T_t) p \dx \\
 & & + \int_{\Omega} \bigl( g'(u)( \hat u_t -u) - g(\hat u_t) + g(u)) p \dx=: \sum_{k=1}^4 \tilde T_k.
 \end{eqnarray*}
 Arguing as in \cite[(3.14)]{DHH25-convergence} one shows that 
\begin{equation} \label{eq:tildet1} 
\tilde T_1  \leq t \int_{\Omega} \Div( f V) p  \dx +  c t^2 + c t  \sup_{0 \leq \sigma \leq t} \Vert \nabla f \circ T_\sigma - \nabla f \Vert_{L^2}.  
\end{equation}
 Since  $A_t = (DT_t)^{-1} (DT_t)^{-\mathsf{T}} \mbox{det} DT_t$ and $| DV| \leq 1$
 one verifies that
\begin{equation} \label{idmat}
I - A_t = t \bigl( D V + D V^{\mathsf{T}} - \Div V I  \bigr) + R_2, \qquad \mbox{ with } | R_2  | \leq c t^2,
\end{equation}
where $c$ only depends on $d$. Therefore
\begin{eqnarray} 
\tilde T_2 & = & \int_{\Omega} (I- A_t) \nabla p \cdot \nabla u  \dx + \int_{\Omega} (I - A_t) \nabla p \cdot \nabla ( \hat u_t - u)  \dx \nonumber \\
& = & 
t \int_{\Omega} \bigl( D V + D V^{\mathsf{T}} - \Div V I  \bigr) \nabla u \cdot \nabla p  \dx 
+ \int_{\Omega} R_2 \nabla u \cdot \nabla p  \dx \nonumber \\
& & + \int_{\Omega} (I - A_t) \nabla p \cdot \nabla ( \hat u_t - u)  \dx 
\nonumber\\
& \leq &  t \int_{\Omega} \bigl( D V + D V^{\mathsf{T}} - \Div V I  \bigr) \nabla u \cdot \nabla p  \dx + c \Vert \nabla p \Vert_{L^2} \bigl( t^2 +  t \Vert \hat u_t - u \Vert_{H^1} \bigr),  \label{eq:tildet2}
\end{eqnarray}
where we also used  \eqref{discapriori}. In order to deal with $\tilde T_3$ we use \eqref{dif1} and \eqref{discapriori} to obtain
\begin{eqnarray}
\tilde T_3 & =& \int_{\Omega} g(u) p (1 - \mbox{det} DT_t) \dx + \int_{\Omega} (g(\hat u_t - g(u)) p (1 - \mbox{det} D T_t) \dx \nonumber \\
& \leq & -t\int_{\Omega} g(u)p \mbox{div}V \dx + c t^2+ c t  \Vert \hat u_t - u \Vert_{L^2} \Vert p \Vert_{L^2} \nonumber \\
& \leq &  -t\int_{\Omega} g(u)p \mbox{div}V \dx + c t^2+ c \Vert \hat u_t - u \Vert_{L^2}^2. \label{eq:tilde3}
\end{eqnarray}
Finally, Taylor expansion yields
\[
g(\hat u_t)-g(u) = g'(u)(\hat u_t-u) + \int\limits_0^1 (1-s)g''(u+s(\hat u_t-u))ds (\hat u_t-u)^2
\]
and hence
\begin{eqnarray}
\tilde T_4 & = & 
- \int_{\Omega} p \int\limits_0^1 (1-s)g''(u+s(\hat u_t-u))ds (\hat u_t-u)^2 \dx \label{eq:tilde4} \\
& \leq & c 
 \int_{\Omega} \vert \hat u_t-u\vert^2  \vert p\vert \dx \le c 
 \Vert \hat u_t -u\Vert_{L^4}^2  \Vert p\Vert_{L^2} \leq  c \Vert \hat u_t-u\Vert_{H^1}^2. \nonumber
\end{eqnarray}
Collecting the above terms we have
\begin{eqnarray}
T_{2,4} & \leq &  t \int_{\Omega} \bigl( D V + D V^{\mathsf{T}} - \Div V I  \bigr) \nabla u \cdot \nabla p  \dx + 
t \int_{\Omega} \Div( f V) p  \dx  \nonumber \\
& & +  c t^2 +c \Vert \hat u_t - u \Vert_{H^1}^2 + c t  \sup_{0 \leq \sigma \leq t} \Vert \nabla f \circ T_\sigma - \nabla f \Vert_{L^2} -t\int_{\Omega} g(u)p \mbox{div}Vdx. \label{t24}
\end{eqnarray}

In conclusion,
\begin{eqnarray}
T_2 & \leq &  t  \int_{\Omega} j_x(\cdot,u,\nabla u) \cdot V  \dx - t \int_{\Omega} j_z(\cdot,u,\nabla u) \cdot DV^{\mathsf{T}} \nabla u \nonumber \\
& & +  t \int_{\Omega} \bigl( D V + D V^{\mathsf{T}} - \Div V I  \bigr) \nabla u \cdot \nabla p  \dx + 
t \int_{\Omega} \Div( f V) p -t\int_{\Omega} g(u)p \mbox{div}V  \dx  \nonumber \\
& & +  c t^2 +c \Vert \hat u_t - u \Vert_{H^1}^2 + c t  \sup_{0 \leq \sigma \leq t} \Vert \nabla f \circ T_\sigma - \nabla f \Vert_{L^2}. \label{t2a}
\end{eqnarray}
Finally, $T_3$ can be estimated as in \cite[(3.19)]{DHH25-convergence},
so that 
\begin{equation} \label{t3}
T_3 \leq ct \bigl( t + \Vert \hat u_t - u \Vert_{H^1} \bigr) \leq c t^2 + c \Vert \hat u_t - u \Vert_{H^1}^2.
\end{equation}
If we insert the estimates \eqref{eq:t1}, \eqref{t2a} and \eqref{t3} into \eqref{difj} and recall \eqref{firstvar} we obtain
\begin{equation} \label{difj1}
\mathcal J(\Omega_t) - \mathcal J(\Omega)  \leq    t \mathcal J'(\Omega)[V]  + ct \bigl( t  +  \sup_{0 \leq \sigma \leq t} \Vert \nabla f \circ T_\sigma - \nabla f \Vert_{L^2} \bigr)
+ c \Vert \hat u_t - u \Vert^2_{H^1}.
\end{equation}
It remains to bound $\Vert \hat u_t - u \Vert_{H^1}$. To do so,  we use \eqref{state} and \eqref{stateuht} and derive
\begin{eqnarray*}
\lefteqn{\int_{\Omega} A_t  \nabla ( \hat u_t - u) \cdot \nabla \eta  \dx + \int_{\Omega}\big(g(\hat u_t)-g(u)\big)\eta \dx}\\ &=& \int_{\Omega} (I-A_t) \nabla u \cdot \nabla \eta  \dx + \int_{\Omega}g(\hat u_t)\eta (1-\mbox{ det }DT_t) \dx + \int_{\Omega} \bigl(
f \circ T_t  \, \mbox{det} D T_t -f \bigr) \eta  \dx
\end{eqnarray*}
for all $\eta \in H^1_0(\Omega)$.
It follows from  \eqref{idmat} that there exists $0< \delta_2 \leq \delta_1$ such that $A_t \xi \cdot \xi \geq \frac{1}{2} | \xi |^2$ for all $\xi \in \mathbb R^d$ and $0 \leq t \leq \delta_2$. Inserting $\eta=
\hat u_t-u$ into the above relation and using \eqref{idmat} as well as the
fact that $g' \geq 0$  we obtain
\begin{equation} \label{eq:hatu}
\frac{1}{2} \int_D |  \nabla (\hat u_t - u) |^2  \dx   \leq  c t \Vert \nabla u \Vert_{L^2} \Vert \nabla (\hat u_t - u) \Vert_{L^2}
+ c t(1+ \Vert  f \circ T_t  \, \mbox{det} D T_t -f \Vert_{L^2})  \Vert \hat u_t -u \Vert_{L^2}.
\end{equation}
One easily infers from \cite[(3.13)]{DHH25-convergence} that
\begin{displaymath}
\Vert  f \circ T_t  \, \mbox{det} D T_t -f \Vert_{L^2} \leq c
\Vert f \Vert_{H^1},
\end{displaymath}
so that \eqref{eq:hatu} along with Poincar\'e's inequality yields
$ \Vert \hat u_t - u \Vert_{H^1} \leq ct$.
Using this estimate in 
\eqref{difj1} and applying the assumption  $\mathcal J'(\Omega)[V] \leq - \epsilon$ we obtain
\begin{equation} \label{eq:jdecrease}
\mathcal J(\Omega_t) - \mathcal J(\Omega) 
 \leq  \gamma t \mathcal J'(\Omega)[V] -(1-\gamma) \epsilon t   + c t^2  + c t  \sup_{0 \leq \sigma \leq t} \Vert \nabla f \circ T_\sigma - \nabla f \Vert_{L^2}. 
\end{equation}
Arguing as after (3.21) in \cite{DHH25-convergence} we obtain
the existence of $0< \delta \leq \delta_2$ such that
\begin{displaymath}
\sup_{0 \leq \sigma \leq t} \Vert \nabla f \circ T_\sigma - \nabla f \Vert_{L^2} \leq 
\frac{1}{2c} (1-\gamma)   \epsilon \quad \mbox{ for } 0 \leq t \leq  \delta.
\end{displaymath}
Thus we deduce from \eqref{eq:jdecrease} for $0 \leq t \leq \delta$
\begin{displaymath}
\mathcal J(\Omega_t) - \mathcal J(\Omega)
 \leq   \gamma t \mathcal J'(\Omega)[V] -(1-\gamma) \epsilon t
 + c t \delta + \frac{1}{2} t(1-\gamma) \epsilon \leq 
 \gamma t \mathcal J'(\Omega)[V]
\end{displaymath}
provided one chooses in addition $\delta \leq \frac{1}{2c}(1-\gamma)
\epsilon$. 
\end{proof}

\begin{thebibliography}{99}
\bibitem[ADJ21]{ADJ21} G.~Allaire, C.~Dapogny, and F.~Jouve. \newblock Shape and topology optimization. \newblock In {\em Differential Geometric Partial Differential Equations: Part II}, volume~22 of {\em Handbook of Numerical Analysis}, pages 3--124, Amsterdam, Netherlands, 2021. Elsevier.
\bibitem[DHH25-]{DHH25-convergence} K.~Deckelnick, P.~J. Herbert, and M.~Hinze. \newblock {Convergence of a steepest descent algorithm in shape optimisation using $W^{1,\infty}$ functions}. \newblock {\em ESAIM: M2AN}, 59(3):1505--1529, 2025.
\bibitem[VM25]{VM25} B.~Vexler and D.~Meidner. \newblock {\em {Numerical analysis for elliptic optimal control problems}}, volume~67 of {\em Springer Ser. Comput. Math.} \newblock Springer, Cham, 2025.
\end{thebibliography}
\end{document}
